%%%%%%%%%%%%%%%%%%%%%%%%%%%%%%%%%%%%%%%%%%%%%%%%%%%%%%%%%%%%%%
%% Carlos Segarra's Beamer Presentation Template. All credits
%% to Vincent Labatut from whom I took the template and added
%% my own flavour to it. Kudos to <vincent.labatut@univ-avignon.fr>
%%%%%%%%%%%%%%%%%%%%%%%%%%%%%%%%%%%%%%%%%%%%%%%%%%%%%%%%%%%%%%
% setup Beamer
\documentclass[9pt,    % default is 11pt, use 10pt for more compact slides
%    handout,            % collapse all overlays (=animations) and video-invert console text
    english,            % presentation language (theme supports only french & english)
    xcolor=table,       % colors in the tables
    envcountsect,        % include section number in theorem numbers
    aspectratio=169     % Using 16:9 aspect ratio because 2019
]{beamer}

%%%%%%%%%%%%%%%%%%%%%%%%%%%%%%%%%%%%%%%%%%%%%%%%%%%%%%%%%%%%%%
% setup the theme
%\usepackage{au/sty/beamerthemeAU}         % no option at all
\usepackage[light]{csg-temp/sty/beamerthemeAU}   % the "light" option only changes the title and section pages

%%%%%%%%%%%%%%%%%%%%%%%%%%%%%%%%%%%%%%%%%%%%%%%%%%%%%%%%%%%%%%
% setup side notes
\usepackage{pgfpages}                                   % comment all 3 below lines to hide notes
%\setbeameroption{show notes}                           % alternate content and note slides
%\setbeameroption{show only notes}                      % only note slides
%\setbeameroption{show notes on second screen=right}    % dualscreen: right, left, top, bottom
%\usepackage{enumitem}

%%%%%%%%%%%%%%%%%%%%%%%%%%%%%%%%%%%%%%%%%%%%%%%%%%%%%%%%%%%%%%
% name of the biblatex file
\addbibresource{biblio.bib}

%%%%%%%%%%%%%%%%%%%%%%%%%%%%%%%%%%%%%%%%%%%%%%%%%%%%%%%%%%%%%%
% External Packages
\usepackage{datenumber}
\usepackage{varwidth}

% Math Mode
%\usepackage{algorithm}
%\usepackage{algorithmic}
%\usepackage{algorithm2e}
%\usepackage{multicol}
%\usepackage[noend]{algpseudocode}

%%%%%%%%%%%%%%%%%%%%%%%%%%%%%%%%%%%%%%%%%%%%%%%%%%%%%%%%%%%%%%
% title and subtitle of the presentation (the latter is optional)
\newcommand{\mainTitle}{Checkpoint Restore of Container Meshes}
\newcommand{\secondTitle}{Design, Implementation, and Evaluation}
\subtitle{Decentralized Systems - Project Proposal} % leave empty if no subtitle
\title[\mainTitle] % leave empty for no title in footer
    {\Large \mainTitle \\ \normalsize \secondTitle}
\subtitle{Decentralized Systems - Project Proposal} % leave empty if no subtitle
%\subtitle{Master in Research in Informatics - MIRI}
%%%%%%%%%%%%%%%%%%%%%%%%%%%%%%%%%%%%%%%%%%%%%%%%%%%%%%%%%%%%%%
% date of the presentation (leave empty for no date, default is today)
\date[March 13, 2020] % leave empty for no date in footer
    {Friday March 13, 2020}
    %{\datedayname, \today}
%%%%%%%%%%%%%%%%%%%%%%%%%%%%%%%%%%%%%%%%%%%%%%%%%%%%%%%%%%%%%%
% authors and their affiliations (the latter is optional)
\author[] % leave empty for no author in footer
{Carlos Segarra - \texttt{carlos.segarra@estudiant.upc.edu}}
%{\inst{1} Computer Science Lab, Avignon University -- LIA EA 4128 \texttt{\{firstname.lastname\}@univ-avignon.fr}
%\and \inst{2} Institute of Disruptive Innovation, University of Excellence \texttt{\{firstname.lastname\}@univ-excell.fr}
%}
%%%%%%%%%%%%%%%%%%%%%%%%%%%%%%%%%%%%%%%%%%%%%%%%%%%%%%%%%%%%%%
% optional: additional logo (ex. lab)
%\titlegraphic{\includegraphics[width=3cm,]{images/logo_FME.png}}
% if you want several logos, put them in a box
%\titlegraphic{\parbox{3cm}{\includegraphics[width=3cm,]{images/ceri_logo.pdf}\newline\includegraphics[width=3cm,]{images/lia_logo.pdf}}}
%%%%%%%%%%%%%%%%%%%%%%%%%%%%%%%%%%%%%%%%%%%%%%%%%%%%%%%%%%%%%%

%%%%%%%%%%%%%%%%%%%%%%%%%%%%%%%%%%%%%%%%%%%%%%%%%%%%%%%%%%%%%
% Presentation speciphic packages
% \usepackage{multicol}
% \usepackage[titles]{tocloft}
% \renewcommand{\cftchapfont}{\normalfont\bfseries}
\usetikzlibrary{decorations.pathmorphing, patterns}
\usepackage{tabularx}
\newcolumntype{L}[1]{>{\raggedright\arraybackslash}p{#1}}
\newcolumntype{C}[1]{>{\centering\arraybackslash}p{#1}}
\newcolumntype{R}[1]{>{\raggedleft\arraybackslash}p{#1}}
%%%%%%%%%%%%%%%%%%%%%%%%%%%%%%%%%%%%%%%%%%%%%%%%%%%%%%%%%%%%%

%%%%%%%%%%%%%%%%%%%%%%%%%%%%%%%%%%%%%%%%%%%%%%%%%%%%%%%%%%%%%%
\begin{document}
%%% title page
\begin{frame}
  \titlepage
\end{frame}

\begin{frame}
    \frametitle{Project Outline}
    \framesubtitle{Motivation \& Problem Statement}


    \textbf{\textcolor{blue}{Problem Motivation*:}}
    \begin{itemize}
        \item \textbf{\textcolor{blue}{Containers}} have become the go-to alternative for sandboxing and deploying distributed applications. 
        \item \textbf{\textcolor{blue}{Orchestrators}} such as \href{https://kubernetes.io/}{Kubernetes} or \href{https://docs.docker.com/engine/swarm/}{Docker Swarm} ease these deployments and the network configuration, making it transparent to the user.
        \item \href{https://en.wikipedia.org/wiki/Application_checkpointing}{Checkpointing} provides systems with \textbf{\textcolor{blue}{fault-tolerance}} and \textbf{\textcolor{blue}{state consistency}} enabling rollback and restore from previous stable versions.
        \item \textbf{\textcolor{blue}{Live migration}} consists in transparently relocating running services for improved resource-management and increased QoS.
        \item Support for checkpointing (Docker) containers is very limited (\href{https://github.com/docker/cli/blob/master/experimental/checkpoint-restore.md}{[1]}, \href{https://criu.org/Docker}{[2]}) and only in experimental mode. 
    \end{itemize}

    \vspace{10pt}
    \small
    \begin{description}
        \item *This work is complementary, and builds on, my MSc Thesis under the supervision of Jordi Guitart.
    \end{description}
\end{frame}

\begin{frame}
    \frametitle{Project Outline}
    \framesubtitle{Goals, Experiments, \& Tools}

    \vspace{-25pt}
    \begin{alertblock}{Main Goal}
        Design, implementation, and evaluation of a distributed checkpoint/restore algorithm for a container mesh.
    \end{alertblock}

    \begin{columns}[t]
        \begin{column}{.45\textwidth}
            \textbf{\textcolor{blue}{Milestones:}}
            \begin{enumerate}
                \item Deployment of a distributed container mesh.
                \item Design \& Implementation of a (distributed) DNS.
                \item Implementation of a (set of) distributed snapshot algorithm(s).
                \item Transparent network restore.
            \end{enumerate}
        \end{column} \hfill
        \begin{column}{.45\textwidth}
            \textbf{\textcolor{blue}{Tools:}}
            \begin{itemize}
                \item \href{https://github.com/opencontainers/runc}{\texttt{runc}} and \href{https://containerd.io/}{\texttt{containerd}} as container client and runtime.
                \item \href{https://criu.org/Main_Page}{Checkpoint Restore In Userspace (\texttt{criu})} for local checkpointing of containers.
                \item The implementation will be done in \texttt{C} to leverage \href{https://criu.org/C_API}{\texttt{libcriu}}, and \href{http://www.libssh.org/}{\texttt{libssh}} for the network programming.
            \end{itemize}
        \end{column}
    \end{columns}
\end{frame}

\begin{frame}
    \frametitle{Milestones}
    \framesubtitle{A Detailed Breakdown}
    
    \vspace{-15pt}

        \begin{enumerate}
            \item Deployment of a distributed container mesh.
            \begin{itemize}
                \item \texttt{runc} and \texttt{containerd} are not as integrated with S.o.A orchestrators, so I will implement an \textit{ad-hoc}, simple, deployment script.
                \item I will stick to lower-level containerization tools due to better support of checkpoint-restore.
            \end{itemize}
            \item Design \& Implementation of a (distributed) DNS.
            \begin{itemize}
                \item It will keep track of the static IPs assigned to each container, and the dependencies (connections) among them.
                \item The DNS will handle re-connections.
            \end{itemize}
            \item Implementation of a (set of) distributed snapshot algorithms.
            \begin{itemize}
                \item Initially I will implement the well-known \href{https://en.wikipedia.org/wiki/Chandy-Lamport_algorithm}{Chandy-Lamport algorithm}.
                \item If time allows I might explore other algorithms: 
            \end{itemize}
            \item Transparent network restore.
            \begin{itemize}
                \item Different types of mesh restore will be tested. Namely: whole restore (all containers go down and are restored, maybe elsewhere), and single-container restore.
                \item The latter simulates the situation in which a node temporarily becomes unavailable.
            \end{itemize}
        \end{enumerate}

\end{frame}

\begin{frame}
    \frametitle{Planned Experiments}
    
    \begin{columns}[t]
        \begin{column}{.4\textwidth}
            \begin{enumerate}
                \item[\textbf{E1}] \textbf{Microbenchmarks:}
                \begin{itemize}
                    \item DNS Evaluation: scalability \& liveness.
                    \item Distributed Checkpoint Evaluation: elapsed time.
                \end{itemize}
                \item[\textbf{E2}] \textbf{Macrobenchmarks:}
                \begin{itemize}
                    \item Single container fault+restore: single and following pattern.
                    \item Whole mesh checkpoint-restore.
                \end{itemize}
                \item[\textbf{E3}] \textbf{Comparison against SoA:}
                \begin{itemize}
                    \item Repeat E2 with Virtual Machines.
                    \item Repeat E2 with Docker Swarm (w/out CR).
                \end{itemize}
            \end{enumerate}
        \end{column} \hfill
        \begin{column}{.6\textwidth}
            \begin{figure}[h!]
                \centering
                \includegraphics[width=\textwidth]{./local_benchmark.pdf}
                \caption{Local Checkpoint-Restore of single \texttt{runc} containers with Criu. Each container runs a Redis DB with an increasing number of keys. I report average and stdev  of 1000 runs. (Plot from my thesis)}
            \end{figure}
        \end{column}
    \end{columns}

\end{frame}

\begin{frame}
    \frametitle{Tentative Roadmap}

    \begin{table}[h!]
        \rowcolors{1}{blue!20}{}
        \begin{tabular}{L{1.5cm}L{1.5cm}L{10cm}}
            \hline
            \rowcolor{blue!40}
            \textbf{Weeks} & \textbf{Milestone} & \textbf{Planned Work} \\[3pt]
            \hline \hline
            1-2 & M1 \& M2 & First design and implementation work. Proof of concept. \\[3pt]
            3-4-5 & M2 \& M3 & Strengthen the design and finalize all implementation. \\[3pt]  
            6-7 & E1 \& E2 \& E3 & Design of the different experiments and evaluation. \\[3pt] 
            8 & P1 & 05/05 Group Work Intermediate Presentation \\[3pt] 
            9 & P & Preparatio of different dissemination materials: plots and poster. \\[3pt] 
            10 & P2 & 22/05 Work Presentation: Poster Session \\[3pt] 
            \hline
        \end{tabular}
    \end{table}

    \begin{center}
        \textbf{\textcolor{blue}{Suggestions \& Recommendations Welcome!}} \\ \href{mailto:carlos.segarra@estudiant.upc.edu}{carlos.segarra@estudiant.upc.edu}
    \end{center}

\end{frame}



\end{document}
